\documentclass[french]{yLectureNote}

\title{Fiche de révision}
\subtitle{Synthèse des notions essentielles}
\author{Paulhenry Saux}
\date{\today}
\yLanguage{Français}

\professor{}

\usepackage{graphicx}
\usepackage[utf8]{inputenc}
\usepackage{geometry}
\usepackage{tikz}
\usepackage{tkz-tab}
\usepackage{tabularx}
\usepackage{awesomebox}
\usepackage{menukeys}
\usepackage{fancyhdr}
\usepackage{blindtext}
\usepackage{hyperref}
\usepackage{caption}
\usepackage{pifont}
\usepackage{array}
\usepackage{lipsum}
\usepackage{yFlatTable}
\usepackage{multicol}

\newcommand{\Lim}[1]{\lim\limits_{\substack{#1}}\:}
\renewcommand{\vec}{\overrightarrow}
\newcommand{\bz}{\overline{z}}
\DeclareMathOperator{\card}{card}
\renewcommand{\vec}{\overrightarrow}
\newcommand{\N}[0]{\mathbb{N}}
\newcommand{\R}[0]{\mathbb{R}}
\newcommand{\C}[0]{\mathbb{C}}
\newcommand{\dd}[0]{\mathrm{d}}
\newcommand{\er}{\vec{e_{\rho}}}
\newcommand{\ep}{\vec{e_{\varphi}}}
\newcommand{\pip}{\pi^+}
\newcommand{\pim}{\pi^-}
\newcommand{\mv}{\mathcal{V}}
\newcommand{\E}{\vec{E}(M)}
\newcommand{\B}{\vec{B}(M)}
\newcommand{\V}{V(M)}
\DeclareMathOperator{\Div}{Div}
\newcommand{\rot}{\vec{\mathrm{rot}}}
\newcommand{\grad}{\vec{\mathrm{grad}}}
\newcommand{\lap}{\nabla^2}

\begin{document}
\chapter{Rappels}

\section{Électrostatique}

\subsection{Champ et force électrostatique}

\begin{theorem}
[Champ électrostatique au point M]
    \[\E = \frac{1}{4\pi\epsilon_0}\frac{q_1}{PM^2}\vec{u}\] avec \(\vec{u}\) le vecteur unitaire dans la direction de PM.

\end{theorem}


\begin{theorem}
[Force électrostatique]
    \[\vec{F} = \frac{1}{4\pi\epsilon_0}\frac{q_1q_2}{PM^2}\vec{u}\] avec \(\vec{u}\) le vecteur unitaire dans la direction de PM.

\end{theorem}


\begin{theorem}
[Distrbution volumique]
    $$\E = \iiint \frac{\rho \dd \tau}{4\pi \epsilon_0 PM^2}\vec{u}$$ avec $\dd \tau$ un élément de volume.

\end{theorem}


\begin{definition}
[Caractéristique d'une ligne de champ]
    $$\vec{\dd l}\wedge \E = \vec{0}$$

\end{definition}

\subsection{Théorèmes de Gauss et Maxwell}

\begin{theorem}
[Théorème de Gauss]
    $$\phi = \iint_{S} \E\cdot \vec{\dd S} = \frac{Q_{int}}{\epsilon_0}$$ avec $\vec{\dd S}$ orienté vers l'extérieur

\end{theorem}

\begin{theorem}
[Forme locale du Théorème de Gauss / Maxwell - Gauss]
    $$\Div\E = \frac{\rho}{\epsilon_0}$$ avec $\rho$ La densité volumique de charge électrique

\end{theorem}

\begin{theorem}
[Maxwell-Faraday]
    $$\rot\vec{E} = \vec{0}$$ quand il n'y a pas de variation du champ magnétique.

\end{theorem}

\subsection{Relations de passage et potentiel}

\begin{theorem}
[Relation de passage]
La composante normale à l'interface de $\E$ est discontinue à la traversée d'une interface chargée. On a en effet $$\vec{n_{ext}}\cdot (\E_{ext}-\E_{int}) = \frac{\sigma}{\epsilon_0}$$

La composante tangentielle est continue :$$\vec{n_{ext}}\wedge (\E_{ext}-\E_{int}) = \vec{0}$$ 

\end{theorem}

\begin{theorem}
[Relation du potentiel]
    $$\int_A^B \E \cdot \vec{\dd l} = V(A)-V(B)$$
    $$ \E = -\grad(V)$$

\end{theorem}

\begin{theorem}
[Équation de Poisson]
   $$-\lap V = \frac{\rho}{\epsilon_0}$$

\end{theorem}

\subsection{Courant et densité de courant}

\begin{definition}
[Densité de courant]
    $$\vec{j} = \rho\vec{v}$$ avec $\vec{v}$ avec $\vec{v}$ la vitesse des charges.
    $$\vec{j} = \gamma \E$$ où $\gamma$ est la conductivité.

\end{definition}

\begin{theorem}
[Intensité du courant]
    $$I = \iint_S \vec{j}\cdot \vec{\dd S}$$

\end{theorem}

\begin{theorem}
[Théorème de Coulomb]
    $$\E_{ext} = \frac{\sigma}{\epsilon_0}\vec{n_{ext}}$$

\end{theorem}

\section{Magnétostatique}

\subsection{Champ magnétique et sources}

\begin{theorem}
$$\B = \frac{\mu_0}{4\pi}\frac{q\vec{v}\wedge \vec{PM}}{\vec{PM^3}}$$

\end{theorem}

\begin{theorem}
[Champ magnétique d'un volume]
    $$\B = \frac{\mu_0}{4\pi}\iint_V \frac{\vec{j}\wedge \vec{PM}}{PM^3}\dd \tau$$

\end{theorem}

\begin{theorem}
[Loi de Biot et Savat]
            $$\B = \frac{\mu_0}{4\pi}\int_C \frac{I\vec{\dd l}\wedge \vec{PM}}{PM^3}$$

\end{theorem}

\subsection{Théorèmes d'Ampère et conservation du flux}

\begin{theorem}
[Théorème d'Ampère]
\[\int_C \vec{B}(M)\cdot \dd \vec{l} = \mu_0 \sum I_{\text{entrées par C}} \]
Les courant sont comptés positivement s'ils sont dans le sens de \(\vec{n}\).

\end{theorem}

\begin{theorem}
[Théorème d'Ampère pour une distribution volumique de courant]
 \[\int_C\B\cdot \dd \vec{l} = \mu_0 \iint \vec{j}\cdot \vec{n}\dd S\]

\end{theorem}

\begin{theorem}
[Forme locale du théorème d'Ampère / Maxwell - Ampère]
\[\rot(\vec{B})=\mu_0 \vec{j}\]

\end{theorem}

\begin{theorem}
[Loi de conservation du flux magnétique]

    $$\phi = \iint_S \B\cdot \vec{\dd S} = \vec{0}$$

\end{theorem}

\begin{theorem}
[Forme locale de la conservation]
    $\Div \B = 0$

\end{theorem}

\subsection{Relations de passage, potentiel vecteur et équations de Poisson}

\begin{theorem}
[Relation de passage pour le champ magnétique]
    \begin{itemize}
    \item La composante normalew de $\B$ est continue lors de la traversée d'une nappe : $$ \vec{n_{ext}}\cdot (\B_{ext}-\B_{int})=0$$
    \item La composante tangentielle est discontinue : $$\vec{n_{ext}}\wedge (\B_{ext}-B_{int}) = \mu_0 \vec{j}$$
\end{itemize}

\end{theorem}

\begin{definition}
[Potentiel vecteur]
    On appelle $\vec{A}$ le potentiel vecteur de $\B$ : $$\B = \rot\vec{A}$$

\end{definition}
\begin{theorem}[Expression du potentiel vecteur pour un fil]
    $$\vec{A}(M) = \frac{\mu_0}{4\pi} I\int \frac{\dd \vec{l}}{PM}$$
\end{theorem}
\begin{theorem}
[Équation de Poisson]
    $$\lap \vec{A} = -\mu_0 \vec{j}$$

\end{theorem}

\subsection{Forces}

\begin{definition}
[Force de Lorenz]
    $$\vec{f} = q(\E+\vec{v}\wedge \B)$$ avec v la vitesse de la particule

\end{definition}

\begin{definition}
[Force de Laplace]
    $$\vec{F} = \int \dd \vec{F} = \int I\vec{\dd l} \wedge \B$$

\end{definition}

\end{document}
